\title{xLSTM: Extended Long Short-Term Memory}

\begin{abstract}
The Long Short-Term Memory (LSTM) architecture was shaped in the 1990s by the introduction of the constant error carousel and gating mechanisms. Over the years, LSTM networks have demonstrated persistent utility in deep learning, significantly contributing to early Large Language Models (LLMs). However, with the rise of Transformer models—powered by scalable, parallel self-attention mechanisms—the prominence of LSTMs diminished at larger scales. This work revisits the potential of LSTMs by asking: To what extent can LSTMs be enhanced for language modeling when scaled to billions of parameters and augmented with contemporary strategies borrowed from modern LLMs, all while addressing their established shortcomings? Our contributions are twofold: First, we propose exponential gating with specialized normalization and stabilization strategies. Second, we reformulate the memory component of LSTMs, resulting in two variants: (i) sLSTM, which uses a scalar memory representation, scalar updates, and a novel memory mixing procedure, and (ii) mLSTM, which achieves full parallelization by employing matrix memory and a covariance-based update. These LSTM improvements are integrated within residual block frameworks, producing xLSTM blocks; stacking these residually forms xLSTM architectures. Both the exponential gating and revised memory approaches enhance xLSTM performance, leading these models to rival or outperform advanced Transformers and State Space Models with respect to performance and scaling.\\
Code available at: \url{https://github.com/NX-AI/xlstm}
\end{abstract}

\section{Introduction}
\label{sec:intro}

The Long Short-Term Memory (LSTM) ideas~\citep{Hochreiter:91,Hochreiter:97nips,Hochreiter:97}, i.e., the constant error carousel and gating, were introduced to overcome the vanishing gradient problem of recurrent neural networks~\citep{Hochreiter:91,Hochreiter:00book}:

\begin{align}
\label{eq:lstm_idea}
  c_t \ = \ f_t \ c_{t-1} \ + \ 
          i_t \  z_t \ , \quad  
          h_t \ = \ o_t \ \psi({c_t}) \ . 
\end{align}

The cell state's constant error carousel operates through an additive modification to $c_{t-1}$ (depicted in green), which is adjusted by the cell inputs $z_t$ and regulated via sigmoid gating mechanisms (indicated in blue). The input gate $i_t$ and forget gate $f_t$ determine the extent of this alteration, whereas the output gate $o_t$ modulates the exposure of the cell's memory as the hidden state $h_t$. Prior to generating the hidden state, the cell state undergoes normalization or squashing through $\psi$, followed by output gating to yield $h_t$.

LSTMs have demonstrated substantial success across a broad range of fields \citep{Hochreiter:01,Hochreiter:07,Schmidhuber:15}, maintaining dominance in text generation applications up until the emergence of Transformers in 2017~\citep{Vaswani:17}. Their superior performance has been established in a variety of tasks involving sequences, such as text generation~\citep{Graves:13,Karpathy:15blog}, handwriting synthesis~\citep{Graves:13}, sequence-to-sequence translation~\citep{Sutskever:14nips}, program evaluation~\citep{Zaremba:14arxiv}, creating captions for images~\citep{Karpathy:15,Hossain:19}, generating source code~\citep{Karpathy:15blog}, modeling rainfall-runoff processes~\citep{Kratzert:18,Kratzert:19}, and developing hydrological models for flood forecasting~\citep{Nearing:24}. Within reinforcement learning, LSTMs outperform other sequence architectures, as evidenced by their integration into models like AlphaStar for StarCraft II~\citep{Vinyals:17short}, OpenAI Five for Dota 2~\citep{OpenAI:19}, and controllers for magnetic confinement in nuclear fusion~\citep{Degrave:22short}. LSTMs are particularly effective at capturing high-level abstractions, enabling them to extract and retain meaningful semantic features within their memory cells~\citep{Karpathy:15blog}, a property highlighted by discoveries such as neurons encoding numbers and syntax~\citep{Lakretz:19}, language structure~\citep{Bau:19}, and sentiment~\citep{Radford:17arxiv}. Even today, LSTMs continue to play an important role in critical domains~\citep{Degrave:22short,Nearing:24}, evidencing their lasting utility and resilience.

Although LSTMs have achieved significant advances, they are constrained by three primary drawbacks: (i) Inability to update storage choices post hoc. This shortcoming is illustrated with the {\it Nearest Neighbor Search} problem (see also Appendix~\ref{sec:appNearestNeighborSearch}). Here, given a reference vector, the model must traverse a sequence to identify the most similar vector and return the value associated with it at the sequence’s end. The mean squared error for this task is depicted in the left panel of Figure~\ref{fig:lstmProblems}. Standard LSTM architectures find it difficult to replace an already stored value if a better-matching vector is encountered later in the sequence; our proposed xLSTM, utilizing exponential gating, effectively overcomes this issue. (ii) Constraints on storage capacity, as all stored information is reduced to scalar cell states. This is evident in the {\it Rare Token Prediction} scenario. In the right panel of Figure~\ref{fig:lstmProblems}, token prediction perplexity is shown for various frequency-based partitions on the Wikitext-103 dataset~\citep{Merity:17}. LSTMs underperform on infrequent tokens, reflecting their restricted memory. By contrast, our xLSTM introduces a matrix-based memory that ameliorates this limitation. (iii) The architecture’s inherently sequential computation arises from memory blending—where hidden state transitions depend on previous hidden states—thereby hindering parallelization.

Such challenges inherent to the LSTM framework have contributed to the widespread adoption of Transformers~\citep{Vaswani:17} in the context of language modeling. This raises an important question: What levels of performance in language modeling might be unlocked by addressing these constraints and scaling LSTMs to the magnitudes seen in contemporary Large Language Models?

\section{Extended Long Short-Term Memory}
\label{sec:xLSTM}

In response to the known limitations of LSTM, the Extended Long Short-Term Memory (xLSTM) framework implements two significant innovations over the mechanism described by Equation~\eqref{eq:lstm_idea}. These enhancements, specifically exponential gating and alternative memory structures, lead to the development of two distinct LSTM-based models: (i) sLSTM (detailed in Section~\ref{sec:sLSTM}), which utilizes a scalar-valued memory, scalar update mechanism, and incorporates memory mixing; and (ii) mLSTM (explained in Section~\ref{sec:mLSTM}), characterized by its matrix-form memory and an update procedure grounded in covariance or outer product operations, enabling complete parallelization. The exponential gating approach is foundational to the advances introduced in both sLSTM and mLSTM. In particular, mLSTM eliminates memory mixing---that is, the recurrent hidden-to-hidden connections---in order to support parallel computation. Both sLSTM and mLSTM can generalize to configurations involving multiple memory cells, with sLSTM including mechanisms for memory mixing between cells. Moreover, sLSTM can employ multiple heads, where memory mixing is constrained to interactions between cells within a single head, not across different heads. The implementation of heads in sLSTM, coupled with exponential gating, creates new modalities for memory mixing. In the context of mLSTM, multi-head and multi-cell extensions are functionally identical.

These advanced LSTM mechanisms are incorporated into the residual block structure, forming what are referred to as xLSTM blocks (see Section~\ref{sec:xLSTMblock}). Constructing architectures by residual stacking of such xLSTM blocks yields xLSTM architectures (see Section~\ref{sec:xLSTMarchitecure}). The overall structure of xLSTM, along with its main modules, is presented in Figure~\ref{fig:xlstm_sketch}.

\subsection{Review of the Long Short-Term Memory}
\label{sec:orgLSTM}

The initial conceptualization of the LSTM architecture~\citep{Hochreiter:91,Hochreiter:97nips,Hochreiter:97} featured a scalar memory cell as its core mechanism for computation and storage, specifically designed to circumvent the issue of vanishing gradients~\citep{Hochreiter:91,Hochreiter:00book} using what is known as the constant error carousel—namely, the cell state update process. Three gates manage the information flow within the memory cell: the input gate, output gate, and forget gate, with the latter subsequently incorporated by \citet{Gers:00}.  
The LSTM cell’s dynamics at time step $t$ are governed by the following update rules:

\begin{align}
c_t \ &= \  f_t \ c_{t-1} \ + \ i_t \ z_t &  & &\text{cell state} \\
h_t  \ &= \ o_t \ \tilde{h}_t \ , & \tilde{h}_t \ &= \ \psi \left( c_t\right)
  &\text{hidden state} \\
z_t \ &= \ \varphi \left( \tilde{z}_t \right) \ , 
  &\tilde{z}_t \ &=  \ \mathbf{w}^\top_{z} \ \mathbf{x}_t \ + \
  r_{z}  h_{t-1} \ + \  b_{z} \ \
  &\text{cell input} \\
i_t \ &= \ \sigma \left( \tilde{i}_t  \right) \ , 
  &\tilde{i}_t \ &= \ \mathbf{w}^\top_{i} \ \mathbf{x}_t \ + \
  r_{i}  \ h_{t-1} \ + \  b_{i} \ \
  &\text{input gate} \\
f_t \ &= \ \sigma \left(  \tilde{f}_t \right) \ , 
  &\tilde{f}_t \ &= \ \mathbf{w}^\top_{f} \ \mathbf{x}_t  \ + \
  r_{f}  \ h_{t-1} \ + \  b_{f} \ \
  &\text{forget gate} \\
o_t \ &= \ \sigma \left( \tilde{o}_t \right) \ , 
  &\tilde{o}_t  \ &= \ \mathbf{w}^\top_{o} \ \mathbf{x}_t \ + \
  r_{o}  \ h_{t-1} \ + \  b_{o} \ \
  &\text{output gate} 
\end{align}

The vectors $\mathbf{w}_{z}$, $\mathbf{w}_{i}$, $\mathbf{w}_{f}$, and $\mathbf{w}_{o}$ denote the input weight vectors that connect the input $\mathbf{x}_t$ to the cell input, input gate, forget gate, and output gate, in that order. Similarly, $r_{z}$, $r_{i}$, $r_{f}$, and $r_{o}$ represent the recurrent weights linking the previous hidden state $h_{t-1}$ to the cell input, input gate, forget gate, and output gate, respectively. The bias parameters $b_{z}$, $b_{i}$, $b_{f}$, and $b_{o}$ are associated with their respective gates and cell input. The functions $\varphi$ and $\psi$ are non-linearities applied to the cell input and hidden state (commonly implemented as $\tanh$). In particular, $\psi$ serves to limit the range of the cell state, preventing it from growing without bound. Activation functions for all gates are sigmoid, defined as $\sigma \left( x \right) = 1/(1+\exp(-x))$. More recent developments combined several scalar memory cells into a vector, $\mathbf{c}_t \in \mathbb{R}^{d}$, enabling the use of recurrent weight matrices $\mathbf{R} \in \mathbb{R}^{d \times d}$ that allow interactions among the different memory cells' outputs~\citep{Greff:15}; refer to Appendix~\ref{sec:apporgLSTM} for a more thorough discussion. Studies focusing on component removal have indicated that each element within the memory cell architecture is essential for effective performance~\citep{Greff:15}.

\subsection{sLSTM}
\label{sec:sLSTM}

To enhance LSTMs so they can reassess their storage choices, we propose the addition of exponential gates (indicated in red), accompanied by normalization and stabilization mechanisms. Specifically, both input and forget gates may utilize exponential activation functions. For normalization, we define a normalizer state, which accumulates the sum of the input gate multiplied by all subsequent forget gates.

The scalar sLSTM forward pass is:

\begin{align}
\label{eq:slstmforward}
c_t \ &= \  f_t \ c_{t-1} \ + \ i_t \ z_t & & 
 &\text{cell state} \\
n_t \ &= \  f_t \ n_{t-1} \ + \ i_t  & & 
 &\text{normalizer state} \\
h_t \ &= \ o_t \ \tilde{h}_t \ ,
  &\tilde{h}_t \ &= \ c_t / n_t
&\text{hidden state} \\
z_t \ &= \ \varphi \left( \tilde{z}_t \right) \ , 
  &\tilde{z}_t \ &=  \ \mathbf{w}^\top_{z} \ \mathbf{x}_t \ + \
  r_{z}  h_{t-1} \ + \  b_{z}
  &\text{cell input} \\
i_t \ &= \ \!\exp \left( \tilde{i}_t  \right) \ , 
  &\tilde{i}_t \ &= \ \mathbf{w}^\top_{i} \ \mathbf{x}_t \ + \
  r_{i}  \ h_{t-1} \ + \  b_{i}
  &\text{input gate} \\
f_t \ &= \ \sigma \left(  \tilde{f}_t \right) \ \text{OR} \
\!\exp \left(  \tilde{f}_t \right) \ ,
  &\tilde{f}_t \ &= \ \mathbf{w}^\top_{f} \ \mathbf{x}_t  \ + \
  r_{f}  \ h_{t-1} \ + \  b_{f}
  &\text{forget gate} \\
o_t \ &= \ \sigma \left( \tilde{o}_t \right) \ , 
  &\tilde{o}_t  \ &= \ \mathbf{w}^\top_{o} \ \mathbf{x}_t \ + \
  r_{o}  \ h_{t-1} \ + \  b_{o}
  &\text{output gate} 
\end{align}

We transfer the original LSTM gating techniques, i.e., input- and/or hidden-dependent gating plus bias term, to the new architectures. Exponential activation functions can lead to large values that cause overflows. Therefore, we stabilize gates with an additional state \mbox{$m_t$~\citep{Milakov:18arxiv}}:

\begin{align}
\label{eq:slstmstabil}
m_t \ &= \ \max \left( \log ( f_t ) + m_{t-1} , \log ( i_t ) \right) &\text{stabilizer state} \\
i'_t \ &= \ \!\exp \left( \log \left ( i_t \right) - m_t \right) \ = \ \!\exp \left( \tilde{i}_t - m_t \right) \ \, 
  &\text{stabil. input gate} \\
f'_t \ &= \ \!\exp \left( \log \left( f_t \right) + m_{t-1} - m_t \right) \ \, 
  &\text{stabil. forget gate}
\end{align}

In Appendix~\ref{sec:appsLSTM}, we demonstrate that substituting $f_t$ with $f'_t$ and $i_t$ with $i'_t$ during the forward computation leaves both the overall network output and the gradients of the loss with respect to the parameters unchanged.

\paragraph{New Memory Mixing.} Like standard LSTM, sLSTM supports the use of several memory cells (refer to Appendix~\ref{sec:appsLSTM}). By having multiple memory cells, memory mixing is achieved through recurrent links $\mathbf{R}_{\mathbf{z}}$, $\mathbf{R}_{\mathbf{i}}$, $\mathbf{R}_{\mathbf{f}}$, $\mathbf{R}_{\mathbf{o}}$ that transmit information from the hidden state vector $\mathbf{h}$ to the memory cell input $\mathbf{z}$ and to the gates $\mathbf{i}$, $\mathbf{f}$, $\mathbf{o}$. Exponential gating introduces a distinctive element in the process of memory mixing. The updated sLSTM architecture may feature several heads, enabling memory mixing within individual heads, yet not across them. Incorporating heads and exponential gating in sLSTM gives rise to a novel approach for combining memory.

\subsection{mLSTM}
\label{sec:mLSTM}

In order to improve the memory capacity of LSTMs, we replace the traditional scalar memory cell $c \in \mathbb{R}$ with a matrix $\mathbf{C} \in \mathbb{R}^{d \times d}$. This allows retrieval operations to utilize matrix multiplication. At a given time step $t$, a pair of vectors—the key $\mathbf{k}_t \in \mathbb{R}^d$ and the value $\mathbf{v}_t \in \mathbb{R}^d$ (following Transformer nomenclature)—is stored. Subsequently, at time $t + \tau$, the stored value $\mathbf{v}_t$ can be accessed using a query vector $\mathbf{q}_{t+\tau} \in \mathbb{R}^d$. This corresponds to the framework established by Bidirectional Associative Memories (BAMs) \citep{Kohonen:72,Anderson:72,Nakano:72,Anderson:77}.

The covariance update rule~\citep{Sejnowski:77,Dayan:91} for storing a key-value pair is

\begin{align}
\mathbf{C}_t \ &= \ \mathbf{C}_{t-1} \ + \ \mathbf{v}_{t} \ \mathbf{k}_t^\top \ .
\end{align}

In our approach, a layer-norm is applied prior to mapping the inputs to the key and value spaces, ensuring that these representations possess zero mean. The covariance update procedure has been shown to yield optimal results~\citep{Dayan:91} in terms of maximizing the separability among retrieved binary codes, or, equivalently, in optimizing the signal-to-noise ratio. Enhanced separation can be attained under constraints to pairwise associations; however, this entails the acceptance of quadratic computational cost, similar to what is observed in attention mechanisms~\citep{Krotov:16,Krotov:17,Ramsauer:21}. Moreover, the covariance update mechanism aligns with the operation of Fast Weight Programmers \citep{Schmidhuber:92ncfastweights,Schlag:21}, which have subsequently incorporated a fixed decay factor applied to $\mathbf{C}_{t-1}$, as well as a constant learning rate modulating the update $\mathbf{v}_{t} \mathbf{k}_t ^\top$~\citep{Ba:16}. Building on these developments, we incorporate the covariance update rule within the architecture of the LSTM, interpreting the forget gate as controlling the decay, the input gate as determining the learning rate, and utilizing the output gate to modulate the retrieved output vector.

For this matrix memory, the normalizer state is the weighted sum of key vectors, where each key vector is weighted by the input gate and all future forget gates. Again, the normalizer state keeps record of the strength of the gates. Since the dot product between query and normalizer state can be close to zero, we use the absolute value of this dot product and lower bound it by a threshold (typically 1.0) as done previously~\citep{Sun:23arxiv}. The mLSTM forward pass is:

\begin{align}
\label{eq:mlstm_recurrent_begin}
\mathbf{C}_t \ &= \  f_t \ \mathbf{C}_{t-1} \ + \ 
  i_t \ \mathbf{v}_t \ \mathbf{k}_t^\top & &  &\text{cell state} \\
\mathbf{n}_t \ &= \  f_t \ \mathbf{n}_{t-1} \ + \ 
  i_t \ \mathbf{k}_t & &  &\text{normalizer state} \\
\mathbf{h}_t  \ &= \ \mathbf{o}_t \ \odot \ \tilde{\mathbf{h}}_t
\ , \qquad \qquad 
\tilde{\mathbf{h}}_t \ = \   \mathbf{C}_t \mathbf{q}_t \ / \ 
  \max \left\{ |\mathbf{n}_t^\top \mathbf{q}_t|, 1 \right\} 
   &&&\text{hidden state} \\
\mathbf{q}_t \ &= \ \mathbf{W}_q \ \mathbf{x}_t \ + \ \mathbf{b}_q  & &  &\text{query input} \\
\mathbf{k}_t \ &= \ \frac{1}{\sqrt{d}} \mathbf{W}_k \ \mathbf{x}_t \ + \ \mathbf{b}_k  & &  &\text{key input} \\
\mathbf{v}_t \ &= \ \mathbf{W}_v \ \mathbf{x}_t \ + \ \mathbf{b}_v  & &  &\text{value input} \\
i_t \ &= \ \!\exp \!  \! \left( \tilde{i}_t  \right) \ , 
 \qquad \qquad \ \ \ \ \,
  \tilde{i}_t \ = \ \mathbf{w}^\top_{i} \ \mathbf{x}_t \ + \  b_{i} \ \ \ 
  &&&\text{input gate} \\
f_t \ &= \ \sigma \!  \left(  \tilde{f}_t \right) \ \text{OR} \ \!\exp \!  \! \left(  \tilde{f}_t \right) , \ \ \: \!
\tilde{f}_t \ = \ \mathbf{w}^\top_{f} \ \mathbf{x}_t  \ + \
  b_{f}
  &&& \text{forget gate} \\
\label{eq:mlstm_recurrent_end}
\mathbf{o}_t \ &= \ \sigma \left( \tilde{\mathbf{o}}_t \right) \ , \qquad \qquad \quad \ \ \ \,
  \tilde{\mathbf{o}}_t  \ = \ \mathbf{W}_{\mathbf{o}} \ \mathbf{x}_t \ + \
  \mathbf{b}_{\mathbf{o}}
  &&&\text{output gate} 
\end{align}

mLSTM is capable of utilizing several memory cells similarly to standard LSTM architectures. For mLSTM specifically, employing multiple heads yields an effect equivalent to having multiple cells, given the absence of memory mixing mechanisms. To ensure stability for the exponential gates within mLSTM, the same stabilization strategies applied to sLSTM are adopted (refer to Equation~\ref{eq:slstmstabil}). The lack of memory mixing in mLSTM also allows its recurrence relations to be re-expressed in a form that supports parallelization. Additional information can be found in Appendix~\ref{sec:appmLSTM}.

\subsection{xLSTM Architecture}
\label{sec:xLSTMarch}

\paragraph{xLSTM Blocks.}
\label{sec:xLSTMblock}
The purpose of an xLSTM block is to capture and non-linearly encode previous information into a rich, high-dimensional representation, enhancing the distinction between different histories or contexts. The ability to separate past histories is fundamental for accurate sequence prediction, such as identifying the next token. This approach is grounded in Cover’s Theorem~\citep{Cover:65}, which suggests that non-linear embedding in higher-dimensional spaces increases the probability that patterns will be linearly separable, compared to the original space. We analyze two types of residual block designs: (i) The post up-projection residual block (akin to Transformer architectures), which first non-linearly summarizes historical information within its original space. Subsequently, it performs a linear transformation to a high-dimensional space, introduces a non-linear activation, and then projects linearly back to the original space; this process is illustrated in the left portion of Figure~\ref{fig:Backbones}
and the third column of Figure~\ref{fig:xlstm_sketch}. 
A comprehensive version is provided in Figure~\ref{fig:appsLSTM_detailed}
located in the appendix. (ii) The pre up-projection residual block (similar to State Space Models), begins by applying a linear map to a high-dimensional space,

It condenses prior information through nonlinear operations within a high-dimensional representation, followed by a linear transformation returning the information to the initial space. When employing an sLSTM within an xLSTM block, the up-projection occurs after the sLSTM component. Conversely, with an mLSTM, the up-projection is positioned before the mLSTM, due to the increased memory capacity available in the expanded dimensional space. Readers can consult the left section of Figure~\ref{fig:Backbones}, the third column of Figure~\ref{fig:xlstm_sketch}, or Figure~\ref{fig:appsLSTM_detailed} in the appendix for a more comprehensive explanation.

\paragraph{xLSTM Architecture. }
\label{sec:xLSTMarchitecure}
The xLSTM network topology is built by stacking multiple component blocks in a residual fashion, drawing on strategies from~\citep{Srivastava:15,He:16}. Our approach utilizes standard pre-LayerNorm~\citep{Ba:16b} residual backbone designs, paralleling the architecture prevalent in current Large Language Models. For visual reference, see the rightmost column in Figure~\ref{fig:xlstm_sketch}.

\subsection{Memory and Speed Considerations}

In contrast to Transformer architectures, xLSTM models feature linear computational requirements and maintain constant memory usage regardless of sequence length. The compressive nature of xLSTM memory makes it particularly advantageous for industrial contexts and deployment on edge devices.

mLSTM utilizes a memory structure that operates without learnable parameters, but relies on costly computations involving $d \times d$ matrices for both storage and updating. This approach prioritizes enhanced memory capability at the expense of increased computational demands. However, since these matrix operations can be efficiently parallelized using GPUs, their impact on overall runtime remains minimal.

Although mLSTM can be parallelized in a manner similar to FlashAttention~\citep{Dao:22,Dao:23} or GLA~\citep{Yang:23arxiv}, sLSTM's design—specifically its memory mixing via hidden-hidden connections—precludes parallelization. To address this, we have implemented an optimized CUDA kernel with memory management at the register level, enabling GPU computations that are generally less than twice as slow as those for mLSTM.

\section{Related Work}
\label{sec:related}

\paragraph{Linear Attention.} Numerous strategies have been proposed to address the quadratic computational cost of the Transformer's attention mechanism, aiming to achieve linear complexity with respect to sequence length. The Synthesizer generates synthetic attention weights independently of token interactions~\citep{Tay:20arxiv}. Linformer employs a low-rank matrix representation of self-attention and provides a linear approximation~\citep{Wang:20arxiv}. Linear Transformer introduces a linearized formulation of the attention operation~\citep{Katharopoulos:20}. Performer approximates the attention softmax in a linear manner by utilizing positive orthogonal random features~\citep{Choromanski:21}. Furthermore, efficient long convolutions serve as substitutes for attention in the Structured Global Convolution (SGConv)~\citep{Li:22} and Hyena Hierarchy~\citep{Poli:23}.

\paragraph{State Space Models.} In recent years, State Space Models (SSMs) have gained significant attention due to their linear scaling with sequence length and their strong results relative to Transformers. The Structured State Space sequence model (S4)~\citep{Gu:21} was among the initial models introduced in this category, and subsequent developments include the Diagonal State Space (DSS) model~\citep{Gupta:22}, Gated State Space (GSS) models~\citep{Mehta:22}, the S5 model~\citep{Smith:22}, Bidirectional Gated SSM (BiGS)~\citep{Wang:22}, H3 model~\citep{Fu:23}, and Mamba~\citep{Gu:24arxiv}.

\paragraph{Recurrent Neural Networks.} Recurrent Neural Networks (RNNs) have been proposed as alternatives to Transformer architectures and attention mechanisms, primarily due to their linear scalability with respect to sequence length. Language modeling experiments with Deep Linear Recurrent Units (LRUs) have produced encouraging outcomes~\citep{Orvieto:23, De:24arxiv}, as have trials involving Hierarchically Gated Linear RNN (HGRN)~\citep{Qin:23} and the updated HGRN2~\citep{Qin:24arxiv}. Among notable RNN-based strategies for large-scale language modeling, RWKV~\citep{Peng:23arxivshort,Peng:24arxivshort} demonstrates results that are comparable to Transformer-based models.

\paragraph{Gating.}
Gating represents a fundamental concept initially introduced in LSTM architectures, which has subsequently been re-envisioned and adopted across a broad spectrum of recent methods. This mechanism finds application in models such as HGRN~\citep{Qin:23}, HGRN2~\citep{Qin:24arxiv}, Gated Linear Attention (GLA)~\citep{Yang:23arxiv}, Gated State Space (GSS) models~\citep{Mehta:22}, Bidirectional Gated SSM (BiGS)~\citep{Wang:22}, Moving Average Equipped Gated Attention (MEGA)~\citep{Ma:22}, RWKV~\citep{Peng:23arxivshort}, and Mamba~\citep{Gu:24arxiv}.

\paragraph{Covariance Update Rule.} 
To improve memory storage, the mLSTM cell integrates a matrix memory employing a covariance-based update rule. Techniques leveraging similar update paradigms include Fast Weight Programmers \citep{Schmidhuber:92ncfastweights,Schlag:21}, RWKV-5 and RWKV-6~\citep{Peng:24arxivshort}, Retention~\citep{Sun:23arxiv}, Linear Transformer~\citep{Katharopoulos:20}, as well as HGRN2~\citep{Qin:24arxiv}.

\paragraph{Most Related.} Among existing models, Retention~\citep{Sun:23arxiv}, RWKV~\citep{Peng:23arxivshort,Peng:24arxivshort}, and HGRN2~\citep{Qin:24arxiv} are conceptually most similar to xLSTM. Each incorporates matrix memory and/or gating mechanisms. Notably, unlike the newly introduced sLSTM, these methods do not permit memory mixing. The ability to mix memories is crucial for tackling state tracking tasks, which enhances the expressiveness of LSTMs compared to State Space Models (SSMs) and Transformers~\citep{Merrill:24,Deletang:23}. Accurate state tracking is essential in scenarios such as code evaluation or following entities over extended narratives.

\paragraph{Residually Stacking Architectures.} The xLSTM architectures, consistent with the prevailing trends in large-scale deep learning, are assembled via residual stacking of modular units~\citep{Srivastava:15,He:16}. 

Such a design principle played a pivotal role in advancing deep convolutional networks~\citep{He:16}, as well as Transformer models~\citep{Vaswani:17}. Transformers have become foundational to Large Language Models (LLMs), serving as the driving technology in models such as GPT-3~\citep{Brown:20short}, ChatGPT~\citep{Schulman:22short}, GPT-4~\citep{Achiam:23gpt4short}, Megatron-LM~\citep{Shoeybi:19}, Gopher~\citep{Rae:21short}, ERNIE 3.0 Titan~\citep{Wang:21arxivshort}, GLaM ~\citep{Du:21glamshort}, Chinese M6~\citep{Lin:21}, multilingual AlexaTM 20B~\citep{Soltan:22}, OPT~\citep{Zhang:22opt}, Chinchilla~\citep{Hoffmann:22short}, BLOOM~\citep{Scao:22arxivshort}, GLM-130B~\citep{Zeng:22arxivshort}, LaMDA~\citep{Thoppilan:22short}, PaLM~\citep{Chowdhery:22short}, Llama~\citep{Touvron:23arxiv}, and Gemini~\citep{GeminiTeam:23arxiv,Reid:24arxiv}.

\section{Experiments}
\label{sec:Exp}

This section presents an empirical assessment of xLSTM, contrasting its performance against established approaches within the realm of language modeling. Section~\ref{sec:ExpSyntethic} is dedicated to examining xLSTM's distinctive properties using artificial tasks. In Section~\ref{sec:ExpComparison}, we analyze the validation perplexity scores from several contemporary language models that have been trained on 15B tokens sourced from SlimPajama~\citep{Soboleva:23}, and also carry out ablation studies targeting xLSTM. Additionally, we explore the scaling characteristics of these models following the frameworks outlined in \citet{Kaplan:20} and \citet{Brown:20short}. 

Section~\ref{sec:ExpLanguage} expands the language modeling evaluation, directly comparing xLSTM to the top-performing models from Section~\ref{sec:ExpComparison}—each trained this time on a more extensive set of 300B tokens from SlimPajama~\citep{Soboleva:23}. The analysis covers several aspects: first, the extrapolation capabilities for handling longer context sequences; second, validation perplexity results alongside downstream task performance~\citep{Sutawika:24}; third, evaluation across the 571 textual domains encompassed by the PALOMA language benchmark dataset~\citep{Magnusson:23arxivshort}; and finally, a reassessment of scaling behaviors now that the models are exposed to 20 times greater training data.

\label{sec:xlstm_blocks_notation}
Throughout all conducted experiments, we denote xLSTM[$a$:$b$] as representing the proportion $a/b$ between xLSTM blocks utilizing mLSTM and those employing sLSTM. To illustrate, xLSTM[7:1] specifies a configuration with eight blocks, where seven blocks are of the mLSTM-based type and one block is sLSTM-based. When the total number of blocks is set to 48, this ratio results in 42 mLSTM-based and 6 sLSTM-based blocks. Additionally, in every experiment, we incorporate pre up-projection blocks for mLSTM and post up-projection blocks for sLSTM.

\subsection{Synthetic Tasks and Long Range Arena}
\label{sec:ExpSyntethic}
Initially, we examine how xLSTM's exponential gating combined with memory mixing performs in the context of formal language tasks~\citep{Deletang:23}. Subsequently, we evaluate the capabilities of xLSTM's matrix memory by applying it to the Multi-Query Associative Recall task~\citep{Arora:23arxiv}. Lastly, we analyze xLSTM's efficiency in handling extended sequences using benchmarks from the Long Range Arena~\citep{Tay:21}.

\paragraph{Evaluation of xLSTM’s Exponential Gating Combined with Memory Mixing.} We assess the efficacy of xLSTM’s novel exponential gating mechanism paired with memory mixing, expected to facilitate handling state tracking challenges~\citep{Merrill:24, Merrill:23}. To support length extrapolation experiments, we adapt and enhance the formal language benchmarks introduced in \citet{Deletang:23} so that training involves sequences of multiple lengths. Comprehensive task descriptions and supplementary results are provided in Appendix~\ref{sec:appExpSynthetic-formalLang}. The study benchmarks xLSTM against various architectures, including Transformers, State Space Models, and classic Recurrent Neural Networks. Each model’s accuracy is measured specifically on tokens pertinent to each formal language task, with scores normalized from 0 (chance) to 1 (ideal performance). We examine two-block implementations for these models, namely: xLSTM[0:1] (using sLSTM only), xLSTM[1:0] (using mLSTM only), xLSTM[1:1], Llama, Mamba, RWKV, Retention, Hyena, standard LSTM, and Transformer blocks containing LSTM (LSTM (Block)). Figure~\ref{fig:formal-main} depicts experimental outcomes. Importantly, architectures such as Transformers and State Space Models lacking memory mixing capabilities (and thus unable to track states) fail at tasks involving regular grammars, like the parity benchmark. These findings are consistent with prior literature, confirming that Transformers and State Space Models are fundamentally less expressive than RNNs for such tasks~\citep{Merrill:24, Merrill:23, Deletang:23}.

\paragraph{Test of xLSTM's Memory Capacities on Associative Recall Tasks.} In this study, we assess the novel matrix memory structure of xLSTM by examining its memory capacity through the Multi-Query Associative Recall task~\citep{Arora:23arxiv}. For each input sequence, a collection of key-value pairs is randomly sampled from a substantial vocabulary, which the model is then required to store and accurately retrieve upon request. To create a more demanding version of the original task, the number of key-value pairs is increased to as many as 256, and the sequence context is lengthened to 2048, thereby expanding the testing spectrum for model memory capabilities. We evaluate a series of 2-block models, including Llama, Mamba, RWKV-5, RWKV-6, xLSTM[1:1], and xLSTM[1:0], with performance measured by the models' accuracy in recalling the memorized pairs. Transformer-based models like Llama, known for a memory scaling exponentially with the coding dimension~\citep{Ramsauer:21}, serve as the benchmark for this experiment. Figure~\ref{fig:mqar-main} displays the evaluation results. xLSTM[1:1] consistently achieves the highest recall accuracy among all non-Transformer models, regardless of model size. Notably, inclusion of the sLSTM block does not degrade memory performance; instead, it appears to enhance memory retention, particularly in the hardest case involving 256 key-value pairs. Further experimental outcomes can be found in Appendix~\ref{sec:appExpSynthetic-mqar}, where additional analysis suggests that xLSTM's expanded memory enables successful extrapolation to much longer contexts than those presented during training.

\paragraph{Test of xLSTM's Long Context Capabilities on Long Range Arena.} In order to evaluate how well xLSTM manages extended input sequences and broad contextual dependencies, we conducted a comparison with alternative approaches on the Long Range Arena~\citep{Tay:21}. xLSTM reliably achieves high results across every task, indicating that the xLSTM model is highly adept at addressing the various challenges associated with lengthy context scenarios.

For more details, see Appendix~\ref{sec:appExpSynthetic-lra}.

\subsection{Method Comparison and Ablation Study}
\label{sec:ExpComparison}

In order to explore the central research question posed in our paper—namely, the performance of our novel LSTM variants when applied at scale in language modeling tasks—we conduct experiments by training xLSTMs, Transformers, State Space Models, and additional techniques on 15B tokens sourced from SlimPajama, employing a consistent auto-regressive framework. Subsequently, we evaluate these models using the validation dataset and carry out ablation experiments focused on the xLSTMs.

\paragraph{Benchmarking xLSTM Against Alternative Approaches.} To conduct a comprehensive evaluation, we train all models on 15B tokens drawn from SlimPajama~\citep{Soboleva:23}. Perplexity scores on the validation set serve as the primary metric for assessment. The comparison spans these approaches: our newly proposed xLSTM, GPT-3 (Transformer)~\citep{Brown:20short}, Llama (Transformer)~\citep{Touvron:23arxiv}, H3 (SSM)~\citep{Fu:23}, Mamba (SSM)~\citep{Gu:24arxiv}, RWKV-4 (RNN)~\citep{Peng:23arxivshort}, RWKV-5 (RNN)~\citep{Peng:24arxivshort}, RWKV-6 (RNN)~\citep{Peng:24arxivshort}, GLA (linear Transformer)~\citep{Yang:23arxiv}, HGRN (RNN)~\citep{Qin:23}, HGRN2 (RNN)~\citep{Qin:24arxiv}, RetNet (linear Transformer)~\citep{Sun:23arxiv}, Hyena (linear Transformer)~\citep{Poli:23}, along with xLSTM[1:0] and xLSTM[7:1] configurations as described in Section~\ref{sec:xlstm_blocks_notation}. 

All models were trained using mixed precision, although for RWKV-5, RWKV-6, GLA, and HGRN2, the mixed precision implementation omitted PyTorch's native automated mixed precision (see Appendix Section~\ref{sec:appGenTrainProc} for details). We classify models into three groups: (a) Transformers, (b) State Space Models (SSMs), and (c) Recurrent Neural Networks (RNNs) as well as linear Transformers. Linear Transformers utilize linear alternatives to traditional attention mechanisms. Each model matches the parameter count of a GPT-3 350M model, which is characterized by an embedding dimension of 1024 and 24 residual blocks; GPT-3 uniquely leverages weight sharing between token and output embeddings, thus containing fewer parameters. 

Results shown in Table~\ref{tab:spaj15b_model_results} demonstrate that xLSTM achieves lower validation perplexity than all baseline models. Further methodological specifics are available in Appendix~\ref{sec:appExpComparison}. Figure~\ref{fig:temp_sclaw15B} plots the scaling trends for this suite of experiments, implying that the superior performance of xLSTM is likely to persist as model size increases.

\paragraph{Ablation Studies.}
The findings presented in Table~\ref{tab:spaj15b_model_results} and Figure~\ref{fig:temp_sclaw15B} indicate that xLSTM delivers strong language modeling performance when trained with 15B tokens from SlimPajama. To isolate the effects contributing to the improvement from LSTM to xLSTM, we incrementally transform the standard LSTM structure into the xLSTM framework. The process begins by embedding LSTM layers within pre-LayerNorm residual architectures. Next, this configuration is expanded to incorporate a post up-projection module. In the final step, exponential gating and matrix memory components are introduced.

The results are shown in Table~\ref{tab:ablstudies} (top). 

The ablation experiments demonstrate that the notable gains in performance can be credited to both the exponential gating technique and the inclusion of matrix memory. Furthermore, recognizing the critical role of gating in RNNs and State Space Models, we conduct ablations on various gating configurations. As shown in Table~\ref{tab:ablstudies} (bottom), our results indicate that allowing each gate to be adaptive and responsive to input yields progressively improved outcomes. Detailed examinations of the specific backbone elements are provided in Appendix~\ref{sec:appAblDetails}.

\subsection{xLSTM as Large Language Model}
\label{sec:ExpLanguage}

This work culminates with extensive language modeling experiments aimed at evaluating xLSTM as a candidate for large language models. Accordingly, we scale up the training corpus, utilizing 300B tokens from SlimPajama—this matches the token volume used in previous studies such as Mamba~\citep{Gu:24arxiv} and Griffin~\citep{De:24arxiv}.

For comparison, we benchmark xLSTM against RWKV-4, Llama, and Mamba, each representing the main architecture categories described in Section~\ref{sec:ExpComparison}. RWKV-4 is chosen to represent RNN-based approaches, as stable training precision for RWKV-5, RWKV-6, and HGRN2 was only established after initiation of the 300B token runs (refer to Appendix~\ref{sec:appGenTrainProc}).

We train models across four scales (125M, 350M, 760M, 1.3B parameters), assess their sequence extrapolation abilities, and measure their performance on held-out validation data. Their effectiveness on downstream tasks is also evaluated, alongside their results on language modeling across 571 distinct text domains from the PALOMA benchmark. Finally, scaling law analyses are conducted to further understand their behavior as model capacity increases.

\paragraph{Sequence Length Extrapolation.}
\label{sec:spaj300B_ctx_extrapolation}
To begin, we examine the ability of large 1.3B-parameter variants of xLSTM, RWKV-4, Llama, and Mamba to generalize to longer sequence lengths. These models are each trained using a context window of 2048 tokens and subsequently assessed at increased context lengths, extending up to 16384 tokens. The outcomes of these experiments are depicted in Figure~\ref{fig:spaj300B_seq_extrapolation}. Notably, the xLSTM architecture sustains low perplexity over these extended sequences, distinguishing itself from the alternative methods.

\paragraph{Validation Perplexity and Downstream Tasks.}
\label{sec:spaj_downstream_eval}
Next, we assess the performance of all evaluated model sizes—xLSTM, RWKV-4, Llama, and Mamba—using the SlimPajama validation suite for next-token prediction, in addition to a set of downstream benchmarks designed to probe common sense reasoning abilities.

The validation set perplexities for each method are shown in the third column of Table~\ref{tab:lmeval}. For all evaluated model sizes, xLSTM[1:0] and xLSTM[7:1] consistently yield the lowest validation set perplexities, marking them as the top-performing models by this measure. The remaining columns of Table~\ref{tab:lmeval} summarize results on various downstream tasks. Across nearly all tasks and model sizes, xLSTM achieves superior performance compared to alternative methods, with the exception that Mamba occasionally surpasses xLSTM on the ARC benchmark. Additional experimental details can be found in Appendix~\ref{sec:appExpLanguage}.

\paragraph{Performance on PALOMA Language Tasks.} Next, we evaluate next token prediction accuracy across all model sizes for xLSTM, RWKV-4, Llama, and Mamba using PALOMA language tasks~\citep{Magnusson:23arxivshort}. The evaluation considers perplexity as a metric for next token prediction across 571 diverse text domains, encompassing sources from nytimes.com to Reddit’s r/depression. 

Table~\ref{tab:ppleval} presents the perplexity results, categorizing them into language modeling benchmarks (first seven columns) and specialized domain benchmarks (last five columns). On these tasks, xLSTM[1:0] outperforms xLSTM[7:1].

Analyzing results across domains, xLSTM[1:0] achieves lower perplexity than Mamba in 568 out of 571 domains (99.5\%), lower than Llama in 486 domains (85.1\%), and outperforms RWKV-4 in 570 domains (99.8\%). Further information can be found in Appendix~\ref{sec:appExpLanguage}.

\paragraph{Scaling Laws.} In the fourth evaluation, we examine the power-law scaling patterns, which facilitate predictions of performance at greater model capacities~\citep{Kaplan:20,Brown:20short}. The scaling trends depicted in Figure~\ref{fig:temp_sclaw300B} reveal that while all examined models display comparable scaling trajectories, their baseline levels differ. RWKV-4 exhibits the poorest results, trailed by Llama and then Mamba. Notably, xLSTM outperforms Mamba, with the degree of improvement analogous to that observed between Mamba and Llama. These scaling dynamics suggest that as model size increases, xLSTM maintains an advantageous position relative to Transformer and State-Space model architectures.

\textbf{Generation Times and Maximal Throughput.} In Figure~\ref{fig:inference_time_generation}, we report the latency associated with text generation for the xLSTM variants at the 1.3B parameter scale, and present their peak throughput in Figure~\ref{fig:inference_time_decoding_throughput} (left). For reference, we include comparably sized implementations of Mamba, Llama, and RWKV from HuggingFace, utilizing a static key-value cache for Llama. At the point when these benchmarks were conducted, both full cache compilation of the Transformer architecture and compilation of the Mamba model using \texttt{torch.compile} were unavailable. Generation time comparisons were carried out with all models configured for batch size 1 and a pre-fill of 16 tokens, a setting that notably benefits the Transformer architecture.

The results in Figure~\ref{fig:inference_time_generation} illustrate that xLSTM, Mamba, and RWKV-4 display linear growth in generation time, whereas Llama exhibits the expected quadratic complexity. To assess decoding throughput, varying batch sizes and pre-fill levels are explored for Llama. As shown in Figure~\ref{fig:inference_time_decoding_throughput} (right), xLSTM supports substantially larger batch sizes relative to Llama, thanks to its fixed memory requirements, resulting in superior throughput performance.

\section{Limitations}

(i) Unlike mLSTM, sLSTM’s memory mixing mechanism restricts parallelizable operations, thereby preventing efficient parallel execution. However, we have implemented an accelerated CUDA kernel for sLSTM, which is currently less than twice as slow as our parallelized mLSTM kernel. (ii) The existing mLSTM CUDA kernels lack optimization, resulting in performance that is roughly four times slower than FlashAttention or the scan operation employed in Mamba. Optimized kernels inspired by FlashAttention could significantly improve speed. (iii) mLSTM’s matrix memory incurs substantial computational demands because it involves manipulating $d \times d$ matrices. Nevertheless, both memory updates and retrieval are parameter-free and can leverage standard matrix algebra for parallelism, thus the practical time overhead from this complexity is relatively minor. (iv) Careful selection of initial values for forget gates is necessary. (v) As matrix memory in mLSTM remains unaffected by sequence length, extending the sequence length can potentially saturate memory for longer contexts. Nonetheless, this does not constitute a limitation for context sizes up to 16k, as evidenced in Section~\ref{sec:spaj300B_ctx_extrapolation}. (vi) Owing to substantial computational costs associated with large-scale language experiments, we did not comprehensively optimize either the model architecture or hyperparameters, particularly for larger xLSTM models. We believe thorough optimization will be essential to fully realize the capabilities of xLSTM.

\section{Conclusion}
Our investigation has partially addressed the fundamental inquiry: How effective is language modeling when LSTM architectures are expanded to billions of parameters? Presently, we can assert that such scaling reaches at least the performance levels of leading methods, such as Transformers and State Space Models. Through the introduction of exponential gating combined with memory mixing and a redesigned memory structure, we have advanced the LSTM architecture to xLSTM. Empirical results demonstrate that xLSTM consistently rivals or outperforms Transformer and State Space Model-based language modeling approaches. According to observed scaling laws, further increases in xLSTM model size are likely to establish xLSTM as a robust competitor to prevailing Large Language Models reliant on Transformer frameworks. Beyond language modeling, xLSTM shows promise for broader application in domains like Reinforcement Learning, Time Series Prediction, and physical system modeling.

\end{document}